\documentclass[11pt,a4paper]{article}
\usepackage[margin=2.2cm]{geometry}
\usepackage{amsmath}
\usepackage{booktabs}
\usepackage{array}
\usepackage{hyperref}

\title{NeuronSeqPerform Sonic Implementation Notes}
\author{Project Technical Documentation}
\date{\today}

\begin{document}
\maketitle

\section{Purpose}
This document explains how audio is produced in NeuronSeqPerform from a sonic implementation perspective. It covers:
\begin{itemize}
    \item timing and control-rate logic,
    \item LIF network activity generation and spike pooling,
    \item FM synthesis voice architecture and sample rendering,
    \item parameter modulation via MIDI (MPD218),
    \item practical behavior and constraints relevant to sound.
\end{itemize}

\section{Signal-Flow Overview}
Implementation files of interest:
\begin{itemize}
    \item main loop: \texttt{main.py}
    \item LIF wrapper: \texttt{model/lif_network.py}
    \item synth wrapper: \texttt{model/fm_synth.py}
    \item native FM core: \texttt{synth\_native/FMSynthMetal.mm}
    \item FM preset data: \texttt{model/presets.py}
    \item MIDI mapping: \texttt{control/midi\_handler.py}
\end{itemize}

Per sequencer tick, the logic is:
\begin{enumerate}
    \item Run $N_{\mathrm{lif}}$ micro-steps of LIF simulation.
    \item Accumulate spike states over those micro-steps.
    \item Downmix dynamic LIF grid to fixed synth grid $8 \times 16$.
    \item Trigger FM envelopes for current sequencer column.
    \item Audio callback renders buffer using current FM synth state.
\end{enumerate}

\section{Clocking and Control Rate}
\subsection{Step Duration}
The sequencer target step duration is a 16th-note at current tempo:
\begin{equation}
T_{\mathrm{step}} = \frac{60}{\mathrm{BPM}} \cdot \frac{1}{4}
\end{equation}

Additional timing terms:
\begin{itemize}
    \item swing offset for odd steps,
    \item random jitter scaled by quantization strength.
\end{itemize}

If quantization strength is $q \in [0,1]$, jitter amplitude is proportional to $(1-q)$.

\subsection{Audio Buffering}
Default audio settings from \texttt{config.py}:
\begin{itemize}
    \item sample rate: $44100$ Hz,
    \item block size: $1024$ samples,
    \item stereo output by duplicating mono synth signal to both channels.
\end{itemize}

\section{LIF Network as Sonic Driver}
\subsection{Purpose in this system}
LIF activity is not directly audible. It controls FM operator envelopes via spike events.

\subsection{Topologies and size}
Supported topologies (native and fallback):
\begin{itemize}
    \item Ring,
    \item FullyConnected,
    \item Feedforward,
    \item SparseRandom,
    \item SmallWorld.
\end{itemize}

Neuron count is runtime-modifiable (typically stepped among $\{128,256,512,1024,2048,4096\}$).

\subsection{Membrane update (fallback math)}
The fallback LIF update follows a standard leaky form:
\begin{equation}
\Delta v = \left(-(v-v_{\mathrm{rest}}) + I\right)\frac{\Delta t}{\tau}
\end{equation}
with refractory masking and threshold/reset spike behavior.

Native LIF follows equivalent semantics with GPU implementation and explicit topology-driven connectivity.

\subsection{Dynamic-grid to synth-grid mapping}
Because synth rendering expects a fixed grid of operators ($8 \times 16$), dynamic LIF activity is pooled:
\begin{itemize}
    \item columns mapped by nearest-neighbor index selection,
    \item rows pooled by band-wise logical OR (max-pooling for spikes).
\end{itemize}

This keeps musical structure stable while allowing large or changing neuron fields.

\section{FM Synth Architecture}
\subsection{Voice and operator structure}
The synth uses $16$ preset columns (voices) and $8$ operators per voice.
Operators are arranged in fixed carrier/modulator pairs:
\begin{center}
$(0 \leftarrow 1), (2 \leftarrow 3), (4 \leftarrow 5), (6 \leftarrow 7)$
\end{center}

At runtime, only one sequencer column is the active target voice; transitions use gain ramps.

\subsection{Per-operator frequency}
For column $c$ and operator $o$:
\begin{equation}
f_{c,o} = f^{\mathrm{base}}_c \cdot r_{c,o} \cdot s_{\mathrm{ratio}}
\end{equation}
where:
\begin{itemize}
    \item $f^{\mathrm{base}}_c$ is column base frequency (from root+scale mapping),
    \item $r_{c,o}$ is preset ratio,
    \item $s_{\mathrm{ratio}}$ is global ratio scale parameter.
\end{itemize}

\subsection{Spike-gated envelopes}
At tick transition, for current column $c^*$:
\begin{equation}
\mathrm{env}_{o,c^*} \leftarrow \mathrm{spike}_{o,c^*} \in \{0,1\}
\end{equation}
Then per audio buffer, carrier and modulator envelopes decay with different coefficients:
\begin{equation}
\mathrm{env}^{\mathrm{next}}_{\mathrm{car}} = \alpha_{\mathrm{car}}\,\mathrm{env}_{\mathrm{car}},
\quad
\mathrm{env}^{\mathrm{next}}_{\mathrm{mod}} = \alpha_{\mathrm{mod}}\,\mathrm{env}_{\mathrm{mod}}
\end{equation}

\subsection{Per-pair FM sample equation}
For pair $p$ in a voice:
\begin{equation}
x_p[n] = A_p[n] \sin\left(\phi_c[n] + I_p[n]\sin(\phi_m[n])\right)
\end{equation}
with:
\begin{align}
A_p[n] &\sim \text{level ramp from envelope} \\
I_p[n] &\sim \text{mod index ramp} \times s_{\mathrm{mod}}
\end{align}

Voice sample is average of active pairs, then weighted by voice gain ramp.

\subsection{Voice crossfade and summation}
Each of 16 voices has a per-buffer gain ramp from current to target state. Final pre-limiter sample:
\begin{equation}
y[n] = \sum_{v=1}^{16} g_v[n] \cdot \bar{x}_v[n]
\end{equation}

\subsection{Output soft limiting}
Output is soft-clipped and scaled:
\begin{equation}
y_{\mathrm{out}}[n] = \tanh(0.7\,y[n])\cdot V_{\mathrm{master}}
\end{equation}

\section{Preset and Tonal Mapping}
\subsection{Presets}
\texttt{model/presets.py} defines 16 presets. Each preset stores:
\begin{itemize}
    \item ratios (8 values),
    \item levels (8 values),
    \item modulation indices (8 values).
\end{itemize}

\subsection{Root and scale to base frequencies}
Column base frequencies are generated from:
\begin{itemize}
    \item root MIDI note,
    \item selected scale semitone offsets,
    \item degree and octave wrapping over 16 columns.
\end{itemize}

MIDI-to-frequency conversion:
\begin{equation}
f = 440 \cdot 2^{(m-69)/12}
\end{equation}

\section{MIDI Control Surface (MPD218)}
\subsection{Knobs (CC)}
\begin{center}
\begin{tabular}{>{\raggedright\arraybackslash}p{1.4cm}p{1.2cm}p{7.6cm}}
\toprule
Bank & CC & Main sonic effect \\
\midrule
A & 3,9,12,13,14,15 & tempo, master volume, active pairs, envelope decay speed, FM index scale, quantization strength \\
B & 16,17,18,19,20,21 & threshold, tau, weight scale, LIF drive, swing, FM ratio scale \\
C & 22,23,24,25,26,27 & aftertouch target, topology select, neuron-count select, LIF steps/tick, root select, scale select \\
\bottomrule
\end{tabular}
\end{center}

\subsection{Pads}
\begin{itemize}
    \item Bank A: root-note selection + velocity-dependent drive boost.
    \item Bank B: scale selection.
    \item Bank C: network functions (randomize/reset, topology cycling, neuron count stepping, etc.).
\end{itemize}

\subsection{Aftertouch}
Aftertouch is routed to one selected target (CC22 selects route), including threshold, tau, drive, tempo, quantization, swing, master volume, FM index scale, active pairs, or decay speed.

\section{Native vs Fallback Rendering}
\subsection{Native (preferred)}
Native FM synthesis uses Metal compute in \texttt{FMSynthMetal.mm}. The Python layer passes arrays and scalar parameters through pybind.

\subsection{Fallback}
A Python fallback path exists in \texttt{model/fm\_synth.py} conceptually, but production behavior should be validated against your current branch state. In normal operation on macOS with compiled extension, native path is used.

\section{Practical Sonic Notes}
\begin{itemize}
    \item LIF spikes do not generate audio directly; they gate FM envelopes.
    \item Increasing LIF steps/tick increases temporal density of envelope triggers.
    \item \texttt{active\_pairs} controls spectral complexity by enabling 1..4 carrier-mod pairs.
    \item \texttt{mod\_index\_scale} and \texttt{ratio\_scale} are primary timbre shapers.
    \item \texttt{decay\_speed} changes per-buffer envelope decay and therefore articulation.
    \item Dynamic neuron-count changes can occur live; sequencer accumulation logic resets safely on shape changes.
\end{itemize}

\section{Minimal Mental Model}
For quick reasoning about the sound engine:
\begin{enumerate}
    \item LIF generates a binary activity pattern each tick.
    \item Pattern is pooled to $8 \times 16$ operator triggers.
    \item Current sequencer column selects one FM voice.
    \item Triggered operators start at env=1 and decay.
    \item FM pair sums are crossfaded, mixed, soft-limited, and output as mono-to-stereo.
\end{enumerate}

\end{document}
